\subsection{DEFINITION OF A FUNCTION BY MEANS OF A TERM}

C54. Let T, A be two terms and let x, y be distinct letters. Suppose that x does not appear in A and that y does not appear in either T or A. Let R be the relation `` \(x \in A\) and \(y = T\) ''. The relation R has a graph F with respect to the letters x and y. This graph is functional; its first projection is A, and its second projection is the set of objects of the form T for \(x \in A\) (§1, no. 6). For every \(x \in A\) we have \(F(x) = T\) .

Let B be the set of objects of the form T for \(x \in A\) . Then

\[
\boldsymbol {R} \Longrightarrow ((\boldsymbol {x}, \boldsymbol {y}) \in \boldsymbol {A} \times \boldsymbol {B});
\]

since the assembly denoted by \(A \times B\) contains neither x nor y, R has a graph F with respect to the letters x and y (no. 1). It is clear that the relation

\[
(x, y) \in F \text {   and   } (x, y ^ {\prime}) \in F
\]

implies \(y = y'\) , and hence F is a functional graph. The remaining statements are evident.

If C is a set which contains the set B of objects of the form T for \(x \in A\) (where y does not appear in C), then the function \((F, A, C)\) is also denoted by the notation \(x \to T\) ( \(x \in A\) , \(T \in C\) ); the corresponding assembly in formal mathematics contains neither x nor y and does not depend on the choice of the letter y satisfying the above conditions. When the context is sufficiently explicit we shall permit ourselves the notations \(x \to T(x \in A)\) , \((T)_{x \in A}\) , or \(x \to T\) , and sometimes simply T or \((T)\) . *Thus we may speak of ``the function \(x^{3}\) '', if the context indicates clearly that we mean the mapping \(x \to x^{3}\) of the set of complex numbers into itself. *

\minorheading{Examples}

\begin{enumerate}
\def\labelenumi{(\arabic{enumi})}
\item
  If \(f\) is a mapping of \(\mathbf{A}\) into \(\mathbf{B}\) , the function \(f\) is equal to the function \(x \to f(x)\) ( \(x \in \mathbf{A}\) , \(f(x) \in \mathbf{B}\) ), which is written simply as \(x \to f(x)\) or also \((f_x)_{x \in \mathbf{A}}\) (the latter notation is especially associated with the phrase ``family of elements'' instead of ``function'').
\item
  Let \(\mathbf{G}\) be a set of ordered pairs. The functions
\end{enumerate}

\[
z \rightarrow \operatorname{pr} _ {1} z (z \in \mathbf {G}, \operatorname{pr} _ {1} z \in \operatorname{pr} _ {1} \mathbf {G})
\]

and

\[
z \rightarrow \mathrm{pr} _ {2} z \quad (z \in \mathrm{G}, \mathrm{pr} _ {2} z \in \mathrm{pr} _ {2} \mathrm{G})
\]

are called respectively the first and second coordinate functions on G; they are denoted by \(pr_{1}\) and \(pr_{2}\) when there is no risk of confusion.

